% GENERATED FILE. Edit canonical lessons, then run npm run generate:books.
% canonical-source-hash: fnv1a64:b80d317dc9ab1386
% canonical-lessons: GE-C351-teilzeit, GE-C351-vollzeit, GE-C351-schicht, GE-C351-ueberstunde, GE-C351-feierabend

\chapter{Die Arbeitszeit}
\label{ch:ge-die-arbeitszeit}
\hlchaptermodality{351}

\begin{chapteropening}
\textbf{By the end of this chapter:} \emph{I can talk about working hours: part-time and full-time, shifts, overtime and the end of the working day.}

Complete the last lesson of chapter 351: I can talk about working hours: part-time and full-time, shifts, overtime and the end of the working day.
\end{chapteropening}

\section[die Teilzeit]{die Teilzeit — part-time (work)}

\label{lesson:GE-C351-teilzeit}



\begin{quote}
Before the new one: say the German for to employ, then the German for fixed-term.
\end{quote}

\subsection*{You'll want to know: die Teilzeit}
\textbf{die Teilzeit} — \textquotedblleft{}part-time (work)\textquotedblright{}.

\begin{quote}
Seit der Geburt ihres Sohnes arbeitet sie in Teilzeit.
\end{quote}

\begin{quote}
\emph{Since the birth of her son she has been working part-time.}
\end{quote}

\emph{In Teilzeit arbeiten} or \emph{Teilzeit arbeiten}. \emph{Der Teilzeitjob} — the part-time job.

\begin{cousinweb}
\emph{Teilzeit} joins \emph{Teil} \textquotedblleft{}part\textquotedblright{}, Old High German \emph{teil}, Proto-Germanic \emph{dailą} (EtymDB), to \emph{Zeit} \textquotedblleft{}time\textquotedblright{}, Old High German \emph{zīt}, Proto-Germanic \emph{tīdiz} (EtymDB).
\end{cousinweb}

\begin{grammarlens}[title={What you've built}]
One more word for telling, explaining and describing at B1.
\end{grammarlens}

\subsection*{Your turn}
\begin{itemize}
\raggedright
  \item \emph{Say it:} \emph{die Teilzeit}
  \item \emph{Say it:} \emph{die Teilzeit}, once more
  \item \emph{Say it:} a sentence of your own with \emph{die Teilzeit}
  \item \emph{Recall:} say the German for to employ, then the German for fixed-term, then say \emph{die Teilzeit} again
\end{itemize}

\begin{tcolorbox}[breakable,colback=teal!4,colframe=teal!35!black,title={Before you move on}]
What does die Teilzeit mean? (\textquotedblleft{}Part-time (work)\textquotedblright{}.) Say it once more.
\end{tcolorbox}
\section[die Vollzeit]{die Vollzeit — full-time (work)}

\label{lesson:GE-C351-vollzeit}



\begin{quote}
Before the new one: say the German for fixed-term, then the German for part-time.
\end{quote}

\subsection*{You'll want to know: die Vollzeit}
\textbf{die Vollzeit} — \textquotedblleft{}full-time (work)\textquotedblright{}.

\begin{quote}
Mein Vater arbeitet Vollzeit, meine Mutter halbtags.
\end{quote}

\begin{quote}
\emph{My father works full-time, my mother half-days.}
\end{quote}

\emph{Vollzeit arbeiten} — to work full-time. \emph{Die Vollzeitstelle} — the full-time job.

\begin{cousinweb}
\emph{Vollzeit} joins \emph{voll} \textquotedblleft{}full\textquotedblright{}, Old High German \emph{foll}, Proto-Germanic \emph{fullaz} (EtymDB), to \emph{Zeit} \textquotedblleft{}time\textquotedblright{}, Proto-Germanic \emph{tīdiz} (EtymDB).
\end{cousinweb}

\begin{grammarlens}[title={What you've built}]
One more word for telling, explaining and describing at B1.
\end{grammarlens}

\subsection*{Your turn}
\begin{itemize}
\raggedright
  \item \emph{Say it:} \emph{die Vollzeit}
  \item \emph{Say it:} \emph{die Vollzeit}, once more
  \item \emph{Say it:} a sentence of your own with \emph{die Vollzeit}
  \item \emph{Recall:} say the German for fixed-term, then the German for part-time, then say \emph{die Vollzeit} again
\end{itemize}

\begin{tcolorbox}[breakable,colback=teal!4,colframe=teal!35!black,title={Before you move on}]
What does die Vollzeit mean? (\textquotedblleft{}Full-time (work)\textquotedblright{}.) Say it once more.
\end{tcolorbox}
\section[die Schicht]{die Schicht — the shift (at work); the layer}

\label{lesson:GE-C351-schicht}



\begin{quote}
Before the new one: say the German for part-time, then the German for full-time.
\end{quote}

\subsection*{You'll want to know: die Schicht}
\textbf{die Schicht} — \textquotedblleft{}the shift (at work); the layer\textquotedblright{}.

\begin{quote}
Ihre Schicht beginnt um sechs Uhr morgens.
\end{quote}

\begin{quote}
\emph{Her shift starts at six in the morning.}
\end{quote}

Plural: \emph{die Schichten}. \emph{Die Nachtschicht} — the night shift. Also a layer: \emph{eine Schicht Schnee}.

\begin{cousinweb}
\emph{Schicht} is usually traced to Middle Low German \emph{schicht}, a word of the mines for a layer of rock and for a stint of work. EtymDB gives no history for it, but lists English \emph{shift} as a cousin.
\end{cousinweb}

\begin{grammarlens}[title={What you've built}]
One more word for telling, explaining and describing at B1.
\end{grammarlens}

\subsection*{Your turn}
\begin{itemize}
\raggedright
  \item \emph{Say it:} \emph{die Schicht}
  \item \emph{Say it:} \emph{die Schicht}, once more
  \item \emph{Say it:} a sentence of your own with \emph{die Schicht}
  \item \emph{Recall:} say the German for part-time, then the German for full-time, then say \emph{die Schicht} again
\end{itemize}

\begin{tcolorbox}[breakable,colback=teal!4,colframe=teal!35!black,title={Before you move on}]
What does die Schicht mean? (\textquotedblleft{}The shift (at work); the layer\textquotedblright{}.) Say it once more.
\end{tcolorbox}
\section[die Überstunde]{die Überstunde — the hour of overtime}

\label{lesson:GE-C351-ueberstunde}



\begin{quote}
Before the new one: say the German for full-time, then the German for the shift.
\end{quote}

\subsection*{You'll want to know: die Überstunde}
\textbf{die Überstunde} — \textquotedblleft{}the hour of overtime\textquotedblright{}.

\begin{quote}
Letzte Woche habe ich zehn Überstunden gemacht.
\end{quote}

\begin{quote}
\emph{Last week I worked ten hours of overtime.}
\end{quote}

Mostly plural: \emph{Überstunden machen} — to work overtime.

\begin{cousinweb}
\emph{Überstunde} joins \emph{über} \textquotedblleft{}over\textquotedblright{} to \emph{Stunde} \textquotedblleft{}hour\textquotedblright{}, from Old High German \emph{stunta} \textquotedblleft{}a while\textquotedblright{}; EtymDB lists English \emph{stound} \textquotedblleft{}a while\textquotedblright{} as a cousin.
\end{cousinweb}

\begin{grammarlens}[title={What you've built}]
One more word for telling, explaining and describing at B1.
\end{grammarlens}

\subsection*{Your turn}
\begin{itemize}
\raggedright
  \item \emph{Say it:} \emph{die Überstunde}
  \item \emph{Say it:} \emph{die Überstunde}, once more
  \item \emph{Say it:} a sentence of your own with \emph{die Überstunde}
  \item \emph{Recall:} say the German for full-time, then the German for the shift, then say \emph{die Überstunde} again
\end{itemize}

\begin{tcolorbox}[breakable,colback=teal!4,colframe=teal!35!black,title={Before you move on}]
What does die Überstunde mean? (\textquotedblleft{}The hour of overtime\textquotedblright{}.) Say it once more.
\end{tcolorbox}
\section[der Feierabend]{der Feierabend — the end of the working day; the evening off}

\label{lesson:GE-C351-feierabend}



\begin{quote}
Before the new one: say the German for the shift, then the German for the hour of overtime.
\end{quote}

\subsection*{You'll want to know: der Feierabend}
\textbf{der Feierabend} — \textquotedblleft{}the end of the working day; the evening off\textquotedblright{}.

\begin{quote}
Nach Feierabend gehe ich gern schwimmen.
\end{quote}

\begin{quote}
\emph{After work I like to go swimming.}
\end{quote}

\emph{Feierabend machen} — to finish work for the day. \emph{Schönen Feierabend!} — Enjoy your evening!

\begin{cousinweb}
\emph{Feierabend} joins \emph{Feier}, from Old High German \emph{fīra} and Latin \emph{fēriae} \textquotedblleft{}holidays\textquotedblright{} (EtymDB; Lewis \& Short), to \emph{Abend} \textquotedblleft{}evening\textquotedblright{}, inherited from Old High German \emph{āband} (EtymDB). It is usually said to have first meant the evening before a holiday.
\end{cousinweb}

\begin{grammarlens}[title={What you've built}]
One more word for telling, explaining and describing at B1.
\end{grammarlens}

\subsection*{Your turn}
\begin{itemize}
\raggedright
  \item \emph{Say it:} \emph{der Feierabend}
  \item \emph{Say it:} \emph{der Feierabend}, once more
  \item \emph{Say it:} a sentence of your own with \emph{der Feierabend}
  \item \emph{Recall:} say the German for the shift, then the German for the hour of overtime, then say \emph{der Feierabend} again
\end{itemize}

\begin{tcolorbox}[breakable,colback=teal!4,colframe=teal!35!black,title={Before you move on}]
What does der Feierabend mean? (\textquotedblleft{}The end of the working day; the evening off\textquotedblright{}.) Say it once more.
\end{tcolorbox}
